\subsubsection{Raffinement du module et conservation des deux lecteurs}
\label{mg03:refinamiento}

L'identité \eqref{eq:defecto-integrado} conserve le défaut entier
\(1215\) au moyen du résidu \(54\) et du quotient \(129\). Sa
généralisation exige de déclarer séparément le module de lecture,
la taille du support et le nombre de coordonnées. Ces trois opérations
produisent des lois de dénombrement distinctes à partir des mêmes feuillets APP.

Pour \(N\geq2\), nous fixons le représentant positif
\[
 \rho_N(s)=1+((s-1)\bmod N)\in\{1,\ldots,N\},
 \qquad s=\rho_N(s)+NQ_N(s),\quad Q_N(s)\in\mathbb Z.
\]
Sur \(\{1,\ldots,N\}^2\), les lectures complètes sont
\(\rho_N(i+j)\) et \(\rho_N(ij)\). Leurs sommes sont notées
\(S_+^{\mathrm{full}}(N)\) et \(S_\times^{\mathrm{full}}(N)\).

\begin{proposition}[Loi du défaut pour un module entier]
Pour tout \(N\geq2\),
\begin{align}
 S_+^{\mathrm{full}}(N)&=\frac{N^2(N+1)}2,
 \label{mg03:suma-plena}\\
 S_\times^{\mathrm{full}}(N)&=
 \frac N2\left(N^2+\sum_{i=1}^{N}\gcd(i,N)\right),
 \label{mg03:producto-pleno}\\
 \Delta_N^{\mathrm{full}}&=
 \frac N2\left(\sum_{i=1}^{N}\gcd(i,N)-N\right).
 \label{mg03:defecto-pleno}
\end{align}
\end{proposition}

\begin{proof}
Dans une ligne additive, la translation par \(i\) permute les \(N\) classes et leurs représentants positifs. La somme par ligne est \(N(N+1)/2\), ce qui démontre \eqref{mg03:suma-plena}. Pour une ligne multiplicative, soit \(g=\gcd(i,N)\). L'image de \(j\mapsto ij\pmod N\) contient les \(N/g\) classes multiples de \(g\), et chacune a \(g\) préimages. Les représentants positifs sont \(g,2g,\ldots,N\), dont la somme est
\[
 g\frac{(N/g)(N/g+1)}2=\frac{N(N+g)}{2g}.
\]
La multiplicité \(g\) donne \(N(N+g)/2\) par ligne. En sommant sur
\(i\), on obtient \eqref{mg03:producto-pleno} ; sa différence avec
\eqref{mg03:suma-plena} produit \eqref{mg03:defecto-pleno}.
\end{proof}

\begin{corollary}[Modules puissances de nombres premiers]
Si \(N=p^m\), avec \(p\) premier et \(m\geq1\), alors
\begin{equation}
 \sum_{i=1}^{p^m}\gcd(i,p^m)
 =p^m+m(p-1)p^{m-1},\qquad
 \Delta_{p^m}^{\mathrm{full}}
 =\frac{m(p-1)}2p^{2m-1}.
 \label{mg03:primo-potencia}
\end{equation}
\end{corollary}

\begin{proof}
Pour chaque \(0\leq k<m\), il existe
\(p^{m-k}-p^{m-k-1}\) indices de valuation \(p\)-adique
exactement \(k\). Leur plus grand commun diviseur avec \(p^m\) vaut
\(p^k\), et la contribution de cette strate est
\((p-1)p^{m-1}\). Les \(m\) strates apportent
\(m(p-1)p^{m-1}\), tandis que l'indice \(i=p^m\) apporte
\(p^m\). La substitution dans \eqref{mg03:defecto-pleno} donne la
seconde formule.
\end{proof}

La spécialisation triadique est
\(\Delta_{3^m}^{\mathrm{full}}=m3^{2m-1}\). Les premiers modules
montrent simultanément les deux sommes et leur différence :
\[
\begin{array}{c|r|r|r}
 N&S_+^{\mathrm{full}}&S_\times^{\mathrm{full}}&\Delta_N^{\mathrm{full}}\\\hline
 3&18&21&3\\
 9&405&459&54\\
 27&10206&10935&729\\
 81&269001&277749&8748
\end{array}
\]

Le raffinement périodique maintient en revanche \(\rho_9\) et étend
le support à \(\{1,\ldots,9r\}^2\), avec \(r\geq1\) entier. Chaque paire de résidus modulo
neuf a \(r^2\) représentants dans ce support, tant pour
l'addition que pour la multiplication. Par conséquent,
\begin{equation}
 S_+^{\mathrm{per}}(9r;9)=r^2\,405,\quad
 S_\times^{\mathrm{per}}(9r;9)=r^2\,459,\quad
 \Delta^{\mathrm{per}}(9r;9)=r^2\,54.
 \label{mg03:periodico}
\end{equation}
Pour \(r=3\), les sommes valent \(3645\) et \(4131\), et le défaut vaut \(486\). La lecture complète dans \(27\) produit \(729\) parce qu'elle actualise à la fois le module et ses représentants ; la lecture périodique produit \(486\) parce qu'elle conserve le module neuf et multiplie ses incidences par neuf. Dans les deux lectures, les divisions \(i+j=R^++NQ^+\), \(ij=R^\times+NQ^\times\), avec le module effectivement employé, maintiennent l'identité entière par cellule et sa somme. La donnée qui doit accompagner chaque défaut est ainsi fixée.

\subsubsection{Variation dimensionnelle à module neuf fixé}
\label{mg03:dimension}

Une seconde extension maintient le module et augmente le nombre de
facteurs. Pour \(d\geq1\), soient
\[
 \mathcal A_d=\{1,\ldots,9\}^d,\qquad
 \Sigma_d(x)=\rho_9\!\left(\sum_{j=1}^d x_j\right),\qquad
 \Pi_d(x)=\rho_9\!\left(\prod_{j=1}^d x_j\right).
\]

\begin{proposition}[Dénombrement dimensionnel de l'accumulation et du produit]
Les sommes des deux lecteurs sur \(\mathcal A_d\) sont
\begin{equation}
 S_{+,d}=45\,9^{d-1},\qquad
 S_{\times,d}=9^{d+1}-9(d+3)6^{d-1}.
 \label{mg03:ley-dimensional}
\end{equation}
Par conséquent,
\begin{equation}
 \frac{S_{\times,d}-S_{+,d}}{9^d}
 =4-(d+3)\left(\frac23\right)^{d-1}
 \longrightarrow4.
 \label{mg03:limite-dimensional}
\end{equation}
\end{proposition}

\begin{proof}
Les \(d-1\) premières coordonnées étant fixées, la dernière parcourt une fois
tous les résidus de la somme. Chaque représentant positif apparaît
\(9^{d-1}\) fois, et leur somme est \(45\) ; cela prouve la formule
additive.

Pour le produit, on distingue les unités
\(U=\{1,2,4,5,7,8\}\), les multiples de trois de valuation un
\(T=\{3,6\}\), et le représentant nul \(9\). Les \(6^d\)
mots dont tous les facteurs appartiennent à \(U\) distribuent leurs produits
uniformément entre les six unités : lorsqu'on fixe \(d-1\) facteurs,
le dernier agit par une permutation du groupe. Leur contribution est
\(27\,6^{d-1}\).

Les mots ayant exactement un facteur dans \(T\) et les autres dans \(U\) sont au nombre de \(2d6^{d-1}\). Pour chaque position du facteur singulier et chaque choix des unités restantes, les deux choix \(3,6\) produisent une fois chacun des résidus \(3,6\). La contribution est \(9d6^{d-1}\). Les mots restants contiennent un facteur \(9\), ou au moins deux facteurs dans \(T\), et leur produit a pour représentant \(9\). Leur nombre est \(9^d-6^d-2d6^{d-1}\). En conséquence,
\[
 S_{\times,d}=27\,6^{d-1}+9d6^{d-1}
 +9\bigl(9^d-6^d-2d6^{d-1}\bigr),
\]
qui se simplifie en \eqref{mg03:ley-dimensional}. La division par
\(9^d\), suivie de la soustraction de la moyenne additive \(5\), donne
\eqref{mg03:limite-dimensional} ; le terme
\((d+3)(2/3)^{d-1}\) tend vers zéro.
\end{proof}

En dimension deux réapparaissent \(S_{+,2}=405\),
\(S_{\times,2}=459\) et le défaut total \(54\). La limite
\(4\) appartient au défaut moyen lors de l'augmentation de la dimension. Le
défaut total, le quotient conservé et la moyenne dimensionnelle sont des
lecteurs spécifiés de l'arithmétique APP ; TRIT et TPK transportent
ensuite ces données avec leur orientation, leur module et leur lieu de résidence.
